\documentclass[../../main.tex]{subfiles}
\graphicspath{{../../image/}{../../figure/}}
\begin{document}

\section*{习题}\label{section:ch14-习题}
\addcontentsline{toc}{section}{习题}

\textbf{14.1}（$\star\star$）证明式\eqref{align:ch14-mc-approx}定义的 $\overline{f}$ 是无偏估计，换句话说，其右边项的期望等于 $\mathbb{E}[f(\bm{z})]$。

\textbf{14.2}（$\star\star$）证明式\eqref{align:ch14-mc-approx}定义的 $\overline{f}$ 的方差由式\eqref{align:ch14-mc-var}给出。

\textbf{14.3}（$\star$）假设 $z$ 是在区间 $(0,1)$ 上均匀分布的随机变量，使用 $y=h^{-1}(z)$ 变换 $z$，其中 $h(y)$ 由式\eqref{align:ch14-cdf-transform}给出。证明 $y$ 具有分布 $p(y)$。

\textbf{14.4}（$\star\star$）给定一个在区间 $(0,1)$ 上均匀分布的随机变量 $z$，找到一个转换 $y=f(z)$，使得 $y$ 具有式\eqref{align:ch14-cauchy}给出的柯西分布。

\textbf{14.5}（$\star\star$）假设变量 $z_1$ 和 $z_2$ 在单位圆上均匀分布，进行式\eqref{align:ch14-box-muller-r}和式\eqref{align:ch14-box-muller-joint}给出的变量变换。证明 $(y_1,y_2)$ 具有式\eqref{align:ch14-box-muller-joint}所示的联合分布。

\textbf{14.6}（$\star\star$）假设 $\bm{z}$ 是一个具有零均值和单位协方差矩阵的 $D$ 维高斯分布的随机变量，并且假设正定对称矩阵 $\bm{\Sigma}$ 具有楚列斯基分解 $\bm{\Sigma}=\bm{L}\bm{L}^{\mathsf T}$，其中 $\bm{L}$ 是一个下三角矩阵（即主对角线以下的元素非零，主对角线以上的元素为零）。证明变量 $\bm{y}=\bm{\mu}+\bm{L}\bm{z}$ 服从均值为 $\bm{\mu}$ 和协方差为 $\bm{\Sigma}$ 的高斯分布。这提供了一种利用服从零均值、单位方差的一元高斯样本来生成服从一般多元高斯分布样本的技术。

\textbf{14.7}（$\star\star$）证明拒绝采样确实是从期望的分布 $p(z)$ 中抽取样本。假设提议分布是 $q(z)$。证明样本值 $z$ 被接受的概率由 $\widetilde{p}(z)/kq(z)$ 给出，其中 $\widetilde{p}$ 是与 $p(z)$ 成比例的任何未归一化的分布，并且常数 $k$ 设置为能够确保 $kq(z)\geq\widetilde{p}(z)$ 的最小值。注意，抽取值 $z$ 的概率等于从 $q(z)$ 中抽取该值的概率乘以给定已经抽取了该值的情况下接受该值的概率。利用这一点，结合概率的加法和乘法规则，写出关于 $z$ 的归一化分布形式，并证明它等于 $p(z)$。

\textbf{14.8}（$\star$）假设变量 $z$ 在区间 $[0,1]$ 上均匀分布。证明变量 $y=b\tan(\pi(z-1/2))+c$ 具有式\eqref{align:ch14-scaled-cauchy}给出的柯西分布。

\textbf{14.9}（$\star\star$）使用连续性和归一化的要求，确定自适应拒绝采样中包络分布[式\eqref{align:ch14-adaptive-envelope}]的系数 $k_j$。

\textbf{14.10}（$\star\star$）通过使用\cref{subsection:ch14-标准分布}讨论的从单一指数分布中采样的技术，设计一种算法，用来从式\eqref{align:ch14-adaptive-envelope}定义的分段指数分布中采样。

\textbf{14.11}（$\star$）证明式\eqref{align:ch14-rw-1}～式\eqref{align:ch14-rw-3}定义的整数上的简单随机游走具有性质 $\mathbb{E}[z^{(\tau)}]=\mathbb{E}[z^{(\tau-1)}]+1/2$，并因此通过归纳得出 $\mathbb{E}[(z^{(\tau)})^2]=\tau/2$。

\textbf{14.12}（$\star\star$）证明\cref{subsection:ch14-吉布斯采样}讨论的吉布斯采样满足式\eqref{align:ch14-detailed-balance}定义的详细平衡性质。

\textbf{14.13}（$\star$）考虑一个双峰分布，其中两个模式之间被一个概率为零的狭窄区域隔开。讨论标准的吉布斯采样是不是遍历的，并由此判断是否能正确地从此分布中采样。

\end{document}
